\documentclass[11pt,a4paper]{article}
\usepackage[T1]{fontenc}
\usepackage[utf8]{inputenc}
\usepackage[italian,provide=*]{babel}
\usepackage{mathptmx,microtype}
\usepackage[margin=2.4cm]{geometry}
\usepackage{amsmath,booktabs}
\pagestyle{empty}
\setlength{\parskip}{0.4em}\setlength{\parindent}{0pt}
\begin{document}
\begin{center}
{\Large\bfseries Premio per il rischio di volatilità e strategie non direzionali\\[2pt] in opzioni sull'indice S\&P 500: evidenza dalle catene 2010--2023}\\[8pt]
Giordano Margotti\\[2pt]{\small Novembre 2026}\\[10pt]
{\bfseries Sintesi}
\end{center}

\textbf{Domanda.} Il lavoro verifica se nei prezzi delle opzioni europee sull'indice S\&P 500 esista un premio per il rischio di volatilità e se le strategie non direzionali di vendita di premio (iron condor, credit spread, strangle e strutture affini) riescano a incassarlo al netto dei costi.

\textbf{Dati e metodo.} Catene di fine giornata delle opzioni SPX dal gennaio 2010 al dicembre 2023: 3.499 sedute e 47.699 coppie data--scadenza. La volatilità implicita è ricavata dai prezzi con il modello di Black (1976) sul forward implicito nella parità put-call; la varianza implicita \emph{model-free} segue la formula del Cboe; la varianza realizzata è misurata fino alla scadenza della stessa opzione. La previsione sotto la misura reale è un GJR-GARCH(1,1) con innovazioni $t$ di Student stimato senza informazioni future. I backtest coprono 434 configurazioni di dieci famiglie, per 1.310.509 operazioni simulate, ai prezzi medi e ai prezzi denaro-lettera.

\textbf{Risultati.} (i) La volatilità \emph{model-free} a 30 giorni supera la realizzata di +3{,}45 punti in media ($t$ 6{,}48); la volatilità \emph{at-the-money} la supera di +0{,}72 punti ($t$ 1{,}36). Il premio risiede nelle ali, non nel centro della superficie. (ii) Il premio \emph{at-the-money} è significativo nel 2010--2014 (+1{,}82 punti, $t$ 4{,}02) e nullo nel 2015--2019 e nel 2020--2023. (iii) La volatilità implicita \emph{at-the-money} a 30 giorni è una previsione non distorta della realizzata (pendenza 0{,}908, test congiunto $p=$ 0{,}409) e ha errore inferiore al GARCH e alla volatilità storica. (iv) Le put con delta 0{,}10 a 30 giorni hanno restituito ai compratori una quota pari a 0{,}44 del premio pagato. (v) Lo Sharpe mediano delle configurazioni è +0{,}39 ai prezzi medi e $-$0{,}40 ai prezzi denaro-lettera; la quota con alfa significativamente positivo rispetto all'indice scende dal 16{,}4\% al 2{,}5\%, valore atteso per solo effetto del caso. (vi) I risultati positivi sono concentrati nelle strutture che vendono put e hanno beta azionario positivo; iron condor e straddle hanno risultato netto nullo o negativo; le regole di uscita anticipata riducono il rendimento perché pagano un secondo differenziale.

\begin{center}\small
\begin{tabular}{lrr}\toprule
 & Prezzi medi & Denaro-lettera \\ \midrule
Sharpe mediano, 434 configurazioni & +0{,}39 & $-$0{,}40 \\
Quota con Sharpe positivo (\%) & 64 & 38 \\
Quota con alfa positivo, $t>2$ (\%) & 16{,}4 & 2{,}5 \\
Put venduta $\Delta$0{,}16, 45 giorni: Sharpe & +0{,}99 & +0{,}93 \\
Put credit spread $\Delta$0{,}20/0{,}05, 30 giorni: Sharpe & +0{,}98 & +0{,}83 \\
Iron condor $\Delta$0{,}16/0{,}05, 30 giorni: Sharpe & +0{,}43 & +0{,}14 \\
\bottomrule\end{tabular}
\end{center}

\textbf{Conclusione.} Il premio esiste come compenso per il rischio di coda ed è coerente con un mercato efficiente sul piano informativo nel centro della superficie. Il differenziale denaro-lettera, l'esposizione azionaria implicita e la selezione a posteriori dei parametri spiegano la quasi totalità dei risultati positivi dei backtest. L'ipotesi che l'incasso sistematico del premio tenda, nel lungo periodo e al netto dei costi, a un risultato nullo non è rifiutata dai dati.

\textbf{Parole chiave:} premio per il rischio di volatilità, opzioni su indice, iron condor, put credit spread, GJR-GARCH, costi di transazione, test multipli.
\end{document}
